% ============================================================
%  Chương 2: Cơ sở lý thuyết
% ============================================================

Chương này trình bày các nền tảng lý thuyết được vận dụng trong đề tài. Mỗi khái niệm không chỉ được mô tả một cách trừu tượng mà còn được phân tích trong mối liên hệ trực tiếp với bài toán phát hiện rò rỉ khí gas, qua đó làm rõ vì sao lựa chọn các công nghệ và mô hình tương ứng là phù hợp với yêu cầu của hệ thống.

\section{Internet of Things và Hệ thống Nhà thông minh}

\subsection{Khái niệm Internet of Things}

Internet of Things (IoT) là một hệ sinh thái công nghệ trong đó các đối tượng vật lý (\textit{things}) được trang bị cảm biến, bộ xử lý và khả năng kết nối mạng, cho phép chúng thu thập, trao đổi và xử lý dữ liệu một cách tự động \cite{nwazor2023}. Đặc điểm cốt lõi của IoT là khả năng \textbf{số hóa môi trường vật lý}: chuyển đổi các đại lượng vật lý (nhiệt độ, độ ẩm, nồng độ khí, áp suất...) thành tín hiệu số có thể được phân tích bằng phần mềm.

Trong bối cảnh nhà thông minh, IoT đóng vai trò là tầng giao tiếp giữa môi trường vật lý và các hệ thống ra quyết định. Một hệ thống phát hiện rò rỉ khí gas dựa trên IoT về bản chất là một \textit{cyber-physical system}: trạng thái vật lý (nồng độ khí trong không gian sống) được liên tục đo đạc, truyền lên các tầng xử lý phía trên, và các quyết định từ phía hệ thống (bật quạt, đóng van, cảnh báo) lại tác động ngược trở lại môi trường vật lý.

\subsection{Kiến trúc phân tầng của hệ thống IoT}

Theo các tài liệu chuẩn về IoT, một hệ thống IoT điển hình được tổ chức thành ba đến bốn tầng chức năng:

\begin{itemize}
  \item \textbf{Tầng nhận thức (Perception Layer)}: Bao gồm các cảm biến và bộ truyền động (actuators). Đây là điểm giao tiếp trực tiếp với môi trường vật lý, có nhiệm vụ chuyển đổi tín hiệu vật lý thành tín hiệu số.

  \item \textbf{Tầng mạng (Network Layer)}: Đảm bảo việc truyền tải dữ liệu từ các thiết bị đầu cuối lên các hệ thống xử lý. Tầng này sử dụng các giao thức như WiFi, Zigbee, LoRaWAN, hay quan trọng nhất trong đề tài là MQTT.

  \item \textbf{Tầng xử lý (Processing Layer)}: Tiếp nhận luồng dữ liệu, thực hiện làm sạch, biến đổi, lưu trữ và phân tích. Trong các hệ thống hiện đại, tầng này thường tích hợp các mô hình học máy để rút trích tri thức từ dữ liệu.

  \item \textbf{Tầng ứng dụng (Application Layer)}: Cung cấp giao diện và dịch vụ đến người dùng cuối thông qua dashboard, ứng dụng di động, chatbot, hay API.
\end{itemize}

Trong đề tài này, chúng tôi mở rộng mô hình ba tầng truyền thống thành \textbf{kiến trúc bốn tầng} bằng cách tách tầng truyền thông và tầng xử lý thành hai tầng độc lập. Quyết định thiết kế này xuất phát từ một quan sát quan trọng: trong các hệ thống streaming hiện đại, lớp message broker (như Kafka) đóng vai trò không chỉ là kênh truyền dẫn mà còn là điểm \textit{decoupling} cho phép nhiều consumer xử lý cùng một dòng dữ liệu một cách độc lập. Việc tách bạch giúp:

\begin{enumerate}
  \item Cho phép thay thế từng thành phần độc lập (ví dụ: chuyển từ Kafka sang Pulsar mà không ảnh hưởng đến tầng xử lý).
  \item Hỗ trợ nhiều mô hình xử lý song song trên cùng một stream.
  \item Đơn giản hóa việc gỡ lỗi nhờ ranh giới trách nhiệm rõ ràng giữa các tầng.
\end{enumerate}

\subsection{Vi điều khiển ESP32}

ESP32 là dòng vi điều khiển 32-bit của Espressif Systems, được lựa chọn rộng rãi trong các dự án IoT nhờ sự cân bằng giữa hiệu năng và chi phí. Cụ thể, ESP32 có bộ xử lý lõi kép Xtensa LX6 với xung nhịp lên đến 240 MHz, 520 KB SRAM tích hợp, và quan trọng nhất là \textbf{tích hợp sẵn WiFi 802.11 b/g/n và Bluetooth} – những giao thức then chốt trong môi trường nhà thông minh.

Đối với bài toán cảm biến khí gas, ESP32 đáp ứng đủ các yêu cầu kỹ thuật:
\begin{itemize}
  \item Có nhiều kênh ADC 12-bit để đọc tín hiệu tương tự từ cảm biến MQ.
  \item Hỗ trợ I2C cho giao tiếp với cảm biến nhiệt ẩm DHT22.
  \item Đủ tài nguyên tính toán để chạy thư viện MQTT client và đóng gói payload JSON.
  \item Tiêu thụ điện năng thấp, phù hợp với các thiết bị cần hoạt động liên tục.
\end{itemize}

\subsection{Cảm biến khí gas và nguyên lý hoạt động}

Dòng cảm biến MQ (MQ-2, MQ-4, MQ-5, MQ-6, MQ-135) là loại cảm biến \textit{metal-oxide semiconductor} (MOS) sử dụng vật liệu nhạy khí Tin Dioxide (SnO$_2$). Nguyên lý hoạt động dựa trên hiện tượng: khi phân tử khí mục tiêu (LPG, methane, propane, ...) tiếp xúc với bề mặt vật liệu được đốt nóng, các phản ứng oxy hóa khử xảy ra làm thay đổi nồng độ electron trên bề mặt, dẫn đến biến thiên điện trở. Sự biến thiên này được đo qua mạch chia áp và chuyển đổi thành điện áp đầu ra tương ứng với nồng độ khí.

Một đặc tính quan trọng cần lưu ý là độ nhạy của cảm biến MQ \textbf{phụ thuộc mạnh vào nhiệt độ và độ ẩm môi trường} \cite{tian2019}. Cùng một nồng độ khí thực tế có thể cho ra các giá trị điện áp khác nhau ở các điều kiện môi trường khác nhau. Đây là một trong những lý do chính giải thích vì sao hệ thống cần thu thập đồng thời nhiệt độ và độ ẩm: không chỉ để báo cáo điều kiện môi trường, mà còn để mô hình học máy có thể \textit{học cách hiệu chỉnh ngầm} nồng độ khí theo điều kiện thực tế.

Về đơn vị đo, nồng độ khí được biểu diễn bằng \textbf{ppm} (parts per million – một phần triệu theo thể tích). Theo tiêu chuẩn an toàn quốc tế, ngưỡng nguy hiểm của LPG (propane/butane) là khoảng 1000 ppm, tương ứng với 10\% LEL (Lower Explosive Limit) – mức nồng độ tối thiểu để hỗn hợp khí có thể cháy. Việc chọn ngưỡng 1000 ppm trong đề tài làm điểm phân loại nguy hiểm/an toàn xuất phát chính từ tiêu chuẩn này.

\section{Giao thức MQTT}

\subsection{Lý do lựa chọn MQTT cho IoT}

MQTT (Message Queuing Telemetry Transport) là giao thức nhắn tin được IBM phát triển từ năm 1999 cho các ứng dụng monitoring đường ống dầu, và sau đó được chuẩn hóa bởi OASIS. MQTT trở thành giao thức gần như mặc định cho IoT nhờ ba đặc tính phù hợp với bài toán đặc thù này:

\begin{enumerate}
  \item \textbf{Tải trọng nhẹ (lightweight)}: Header tối thiểu chỉ 2 byte, phù hợp với các thiết bị có RAM và pin hạn chế.
  \item \textbf{Mô hình bất đối xứng (asymmetric)}: Publisher không cần biết về subscriber và ngược lại, giảm sự phụ thuộc cứng giữa các thành phần.
  \item \textbf{Chịu được mất kết nối}: Cơ chế Last Will Testament (LWT) và Retain Message cho phép xử lý các tình huống thiết bị offline đột ngột – điều thường xảy ra trong môi trường IoT \cite{guo2019}.
\end{enumerate}

\subsection{Kiến trúc Publish-Subscribe}

Khác với mô hình client-server truyền thống nơi client trực tiếp gọi server, mô hình publish-subscribe của MQTT đưa vào một thực thể trung gian là \textbf{broker}. Kiến trúc này gồm ba vai trò:

\begin{itemize}
  \item \textbf{Publisher}: Gửi message lên một \textit{topic} cụ thể mà không cần biết ai sẽ nhận. Trong đề tài, publisher là ESP32 hoặc bộ mô phỏng cảm biến.
  \item \textbf{Broker}: Là nút trung tâm tiếp nhận message, lưu trữ tạm thời và phân phối đến các subscriber phù hợp. Đề tài sử dụng Eclipse Mosquitto – một broker mã nguồn mở phổ biến.
  \item \textbf{Subscriber}: Đăng ký nhận message từ một hoặc nhiều topic. Trong đề tài, subscriber chính là dịch vụ cầu nối MQTT-Kafka.
\end{itemize}

Sự gián tiếp này mang lại \textit{loose coupling} – đặc tính then chốt cho phép hệ thống mở rộng dễ dàng: có thể thêm subscriber mới (ví dụ một service ghi log dự phòng) mà không cần thay đổi bất kỳ thiết bị nào.

\subsection{Quality of Service và đảm bảo phân phối}

Trong môi trường IoT, mạng thường không ổn định, dẫn đến rủi ro mất gói tin. MQTT giải quyết vấn đề này thông qua ba mức Quality of Service (QoS):

\begin{itemize}
  \item \textbf{QoS 0 – At most once}: Message gửi đi không có xác nhận, có thể mất. Phù hợp khi mỗi điểm đo không quá quan trọng và tần suất gửi cao.
  \item \textbf{QoS 1 – At least once}: Đảm bảo message đến đích ít nhất một lần thông qua cơ chế PUBACK. Có thể trùng lặp nhưng không mất.
  \item \textbf{QoS 2 – Exactly once}: Đảm bảo đúng một lần thông qua bắt tay bốn bước, overhead cao nhất.
\end{itemize}

Đề tài lựa chọn \textbf{QoS 1} cho topic \texttt{sensors/gas} vì hai lý do: (i) việc mất dữ liệu cảm biến trong giai đoạn rò rỉ là không chấp nhận được, và (ii) hệ thống xử lý phía sau có cơ chế khử trùng lặp dựa trên timestamp, vì vậy việc nhận message trùng không gây sai lệch kết quả.

\section{Apache Kafka và Stream Processing}

\subsection{Vai trò của message broker trong streaming pipeline}

Trong khi MQTT phù hợp cho lớp truyền tải từ thiết bị, nó không được thiết kế để làm \textit{event log} bền vững với khả năng replay. Đây là lúc Apache Kafka phát huy vai trò: Kafka không chỉ là một message queue mà là một \textbf{distributed commit log} – mỗi message được ghi vào ổ đĩa và có thể được đọc lại nhiều lần.

Đặc tính này có giá trị thực tiễn lớn cho bài toán nghiên cứu:
\begin{itemize}
  \item Khi cần huấn luyện lại mô hình LSTM với cấu hình mới, có thể replay dữ liệu lịch sử từ Kafka thay vì phải chạy lại bộ mô phỏng.
  \item Cho phép thử nghiệm bộ điều khiển mới (ví dụ một thuật toán RL khác) trên cùng dataset đã thu thập trước đó, đảm bảo so sánh công bằng.
  \item Khi consumer (ví dụ Spark) bị lỗi và khởi động lại, nó có thể tiếp tục đọc từ offset cuối cùng đã commit, không bị mất dữ liệu.
\end{itemize}

\subsection{Các khái niệm cốt lõi trong Kafka}

\begin{itemize}
  \item \textbf{Topic}: Kênh logic phân loại message stream theo chủ đề. Mỗi topic được phân thành nhiều \textit{partition} để hỗ trợ xử lý song song.
  \item \textbf{Producer}: Thành phần ghi message vào topic. Trong đề tài, producer là dịch vụ MQTT-Kafka bridge.
  \item \textbf{Consumer / Consumer Group}: Thành phần đọc message. Nhiều consumer trong cùng group sẽ chia phần các partition; nhiều group khác nhau sẽ đọc độc lập cùng dữ liệu.
  \item \textbf{Broker}: Server vật lý lưu trữ partition. Cluster Kafka gồm nhiều broker để chịu lỗi.
  \item \textbf{Offset}: Chỉ số tuần tự của message trong partition. Consumer commit offset để đánh dấu vị trí đã xử lý.
\end{itemize}

Đề tài sử dụng ba topic Kafka tương ứng với ba luồng dữ liệu logic:
\begin{itemize}
  \item \texttt{gas.raw.sensor}: Dữ liệu đo thô từ cảm biến, được bridge từ MQTT.
  \item \texttt{gas.alert.events}: Các sự kiện cảnh báo được sinh ra bởi bộ điều khiển.
  \item \texttt{gas.action.events}: Nhật ký các hành động đã thực thi (ALERT, FAN, VALVE).
\end{itemize}

Việc tách thành ba topic riêng (thay vì một topic chung) giúp các consumer chuyên trách (ví dụ Telegram bot chỉ subscribe topic alert) chỉ phải xử lý dữ liệu liên quan, giảm tải và đơn giản hóa code.

\subsection{Apache Spark Structured Streaming}

Apache Spark Structured Streaming là engine xử lý stream xây dựng trên Spark SQL, áp dụng triết lý \textit{"streaming as continuous queries on a table"}: dữ liệu stream được nhìn nhận như một bảng vô hạn không ngừng được thêm hàng mới, và truy vấn SQL/DataFrame được áp dụng liên tục trên bảng đó. Mô hình này có ưu điểm là cho phép sử dụng lại các thao tác quen thuộc trong batch processing cho stream, giảm đáng kể độ phức tạp khi lập trình.

Spark Structured Streaming xử lý dữ liệu theo cơ chế \textbf{micro-batch}: thay vì xử lý từng event riêng lẻ, Spark gom các event đến trong một khoảng thời gian nhỏ (vài chục đến vài trăm milliseconds) thành một batch và xử lý theo lô. Mô hình này có ưu điểm là throughput cao và cung cấp đảm bảo \textbf{exactly-once semantics} thông qua cơ chế checkpointing.

Trong pipeline của đề tài, Spark thực hiện chuỗi các bước:
\begin{enumerate}
  \item Đọc message từ topic \texttt{gas.raw.sensor}, parse JSON payload.
  \item Duy trì sliding window 60 giây cho mỗi \texttt{device\_id} để cung cấp đầu vào cho LSTM.
  \item Gọi mô hình LSTM dự báo xác suất rủi ro 5 phút.
  \item Gọi bộ điều khiển RL để lựa chọn hành động dựa trên trạng thái hiện tại và xác suất rủi ro.
  \item Ghi kết quả vào InfluxDB và publish event ngược lại Kafka.
\end{enumerate}

\section{Mạng Long Short-Term Memory (LSTM)}

\subsection{Bài toán dự báo chuỗi thời gian}

Trong bài toán phát hiện rò rỉ khí gas, nồng độ khí $g_t$ tại thời điểm $t$ phụ thuộc \textit{không chỉ vào trạng thái hiện tại} của môi trường, mà còn vào \textit{lịch sử} các giá trị trước đó. Một sự kiện rò rỉ chậm có biểu hiện đặc trưng là một xu hướng tăng tuyến tính kéo dài hàng phút; ngược lại, dao động ngẫu nhiên ở mức nền chỉ kéo dài vài giây. Để phân biệt được hai tình huống này, mô hình phải có khả năng \textit{ghi nhớ} thông tin từ quá khứ – đây chính là điểm mà các mạng nơ-ron hồi quy (Recurrent Neural Networks) phát huy thế mạnh.

Bài toán cụ thể trong đề tài có thể được phát biểu như sau: cho chuỗi đo lường $\{(g_{t-59}, T_{t-59}, H_{t-59}), \ldots, (g_t, T_t, H_t)\}$ trong 60 giây gần nhất, dự đoán biến mục tiêu:
\begin{equation}
  y_t = \mathbf{1}\left[\max_{\tau \in (t, t+300]} g_\tau > 1000\right]
\end{equation}
Đây là một bài toán \textbf{phân loại nhị phân trên chuỗi thời gian} với chân trời dự báo 5 phút.

\subsection{Recurrent Neural Networks và vấn đề Vanishing Gradient}

Mạng nơ-ron hồi quy (RNN) là mở rộng của mạng feed-forward để xử lý dữ liệu tuần tự. Tại mỗi bước thời gian $t$, RNN duy trì một trạng thái ẩn $h_t$ tổng hợp thông tin từ toàn bộ chuỗi đầu vào trước đó:
\begin{equation}
  h_t = \tanh(W_h h_{t-1} + W_x x_t + b)
\end{equation}

Tuy lý thuyết RNN có thể học các phụ thuộc dài hạn, trong thực tế nó gặp phải vấn đề \textbf{vanishing gradient}. Khi gradient lan truyền ngược qua nhiều bước thời gian, do tính chất nhân chuỗi của đạo hàm, gradient có thể tiến về 0 (vanishing) hoặc bùng nổ (exploding). Hệ quả là RNN không thể học hiệu quả các mối phụ thuộc kéo dài quá vài chục bước thời gian \cite{gambiroza2020}. Đối với bài toán của đề tài (chuỗi 60 bước), vanishing gradient là một rào cản thực sự nếu sử dụng RNN thuần.

\subsection{Cấu trúc LSTM}

LSTM (Long Short-Term Memory) được Hochreiter và Schmidhuber đề xuất năm 1997 \cite{hochreiter1997} nhằm giải quyết trực tiếp vấn đề vanishing gradient. Ý tưởng cốt lõi là bổ sung một \textbf{cell state} $C_t$ riêng biệt với trạng thái ẩn, đóng vai trò như một "băng tải" cho phép thông tin truyền qua nhiều bước thời gian mà không bị biến đổi đáng kể. Việc cập nhật cell state được điều phối bởi ba cổng (gates):

\begin{align}
  f_t &= \sigma(W_f [h_{t-1}, x_t] + b_f) & \text{(Forget gate)} \\
  i_t &= \sigma(W_i [h_{t-1}, x_t] + b_i) & \text{(Input gate)} \\
  \tilde{C}_t &= \tanh(W_C [h_{t-1}, x_t] + b_C) & \text{(Candidate cell state)} \\
  C_t &= f_t \odot C_{t-1} + i_t \odot \tilde{C}_t & \text{(Cập nhật cell state)} \\
  o_t &= \sigma(W_o [h_{t-1}, x_t] + b_o) & \text{(Output gate)} \\
  h_t &= o_t \odot \tanh(C_t) & \text{(Hidden state)}
\end{align}

Trong các công thức trên, $\sigma$ là hàm sigmoid, $\odot$ là phép nhân Hadamard (element-wise). Mỗi cổng có thể được hiểu như một bộ lọc xác định bao nhiêu thông tin được "đi qua":
\begin{itemize}
  \item \textbf{Forget gate} $f_t$ quyết định bao nhiêu phần của cell state cũ được giữ lại.
  \item \textbf{Input gate} $i_t$ quyết định bao nhiêu thông tin mới được ghi vào cell state.
  \item \textbf{Output gate} $o_t$ quyết định bao nhiêu cell state hiện tại được sử dụng để sinh ra hidden state.
\end{itemize}

Quan trọng là phép cộng $C_t = f_t \odot C_{t-1} + i_t \odot \tilde{C}_t$ trong cập nhật cell state cho phép gradient được truyền ngược qua nhiều bước thời gian với mức suy giảm tối thiểu – đây chính là cơ chế giải quyết vấn đề vanishing gradient.

\subsection{Áp dụng LSTM cho dự báo nồng độ khí gas}

Mô hình LSTM trong đề tài được thiết kế với kiến trúc tối giản nhằm đảm bảo khả năng suy luận thời gian thực:

\begin{itemize}
  \item \textbf{Lớp đầu vào}: Chuỗi 60 bước thời gian, mỗi bước có 3 đặc trưng $(gas\_ppm, temperature, humidity)$ đã chuẩn hóa về thang $[0, 1]$.
  \item \textbf{Lớp LSTM}: 32 đơn vị, không trả về sequence (chỉ lấy đầu ra cuối cùng).
  \item \textbf{Lớp Dropout}: tỷ lệ 0.2 để chống overfitting trên dữ liệu mô phỏng.
  \item \textbf{Lớp Dense}: 16 đơn vị với hàm kích hoạt ReLU.
  \item \textbf{Lớp đầu ra}: 1 đơn vị với hàm kích hoạt sigmoid, cho ra xác suất $P(\text{critical trong 5 phút})$.
\end{itemize}

Sơ đồ khối các lớp của mô hình LSTM được trình bày trong Hình~\ref{fig:lstm_arch}, kèm số tham số ước lượng cho từng lớp và các siêu tham số huấn luyện.

\begin{figure}[H]
  \centering
  \includegraphics[width=0.62\linewidth]{lstm_arch.png}
  \caption{Sơ đồ khối kiến trúc LSTM Forecaster: Input(60,3) $\rightarrow$ LSTM(32) $\rightarrow$ Dropout(0,2) $\rightarrow$ Dense(16, ReLU) $\rightarrow$ Dense(1, Sigmoid). Tổng $\sim$5{,}000 tham số, đầu ra là xác suất $p_{5\min}$.}
  \label{fig:lstm_arch}
\end{figure}

Việc lựa chọn kiến trúc nhỏ (tổng tham số khoảng 5,000) là một quyết định thiết kế có chủ đích. Trong các pipeline streaming, mỗi mili-giây độ trễ đều có giá trị, đặc biệt khi mô hình được gọi cho mỗi mẫu đo lường. Mô hình quá lớn dù có thể đạt độ chính xác cao hơn một chút nhưng sẽ tăng đáng kể chi phí suy luận, làm chậm toàn bộ pipeline. Trong bối cảnh dữ liệu mô phỏng có cấu trúc tương đối rõ ràng, mô hình nhỏ vẫn đủ năng lực biểu diễn.

Hàm mất mát được sử dụng là binary cross-entropy với \textbf{trọng số lớp} (class weights) – một kỹ thuật cần thiết khi dữ liệu bị mất cân bằng. Trong bộ dữ liệu của đề tài, tỷ lệ mẫu dương (cần cảnh báo) chỉ chiếm khoảng 15–25\%, vì vậy nếu không có trọng số, mô hình có xu hướng học chỉ dự đoán nhãn 0 để tối ưu accuracy một cách ngụy biện. Trọng số được tính tự động theo công thức:
\begin{equation}
  w_{\text{pos}} = \frac{1 - p_{\text{pos}}}{p_{\text{pos}}}, \quad w_{\text{neg}} = 1
\end{equation}
trong đó $p_{\text{pos}}$ là tỷ lệ mẫu dương trong tập huấn luyện.

\section{Học tăng cường (Reinforcement Learning)}

\subsection{Khái niệm và động lực sử dụng}

Học tăng cường (Reinforcement Learning – RL) là một paradigm của học máy trong đó một \textbf{tác nhân} (agent) học cách ra quyết định thông qua tương tác trực tiếp với một \textbf{môi trường} (environment). Khác với học có giám sát – nơi mô hình được dạy bằng các cặp (đầu vào, đầu ra đúng), RL không yêu cầu nhãn đúng cho từng tình huống. Thay vào đó, agent thực hiện hành động, quan sát kết quả, nhận tín hiệu \textbf{phần thưởng} (reward) phản ánh chất lượng hành động đó, và dần dần học được \textbf{chính sách} (policy) tối ưu để cực đại hóa tổng phần thưởng tích lũy.

\begin{figure}[H]
  \centering
  \includegraphics[width=0.7\textwidth]{rl_loop.png}
  \caption{Vòng lặp tương tác Agent--Environment trong Reinforcement Learning.}
  \label{fig:rl_loop}
\end{figure}

Việc lựa chọn RL cho bài toán điều khiển khí gas xuất phát từ ba quan sát then chốt:

\begin{enumerate}
  \item \textbf{Không tồn tại nhãn "đúng" tuyệt đối}: Trong cùng một trạng thái nồng độ khí 600 ppm với xu hướng tăng, hành động tối ưu có thể là ALERT, FAN hay VALVE tùy thuộc vào lịch sử và tình huống cụ thể. Không có một quy tắc cứng nào có thể bao quát mọi trường hợp.

  \item \textbf{Tín hiệu phản hồi có độ trễ}: Một hành động đóng van ở thời điểm $t$ chỉ có thể đánh giá là tốt hay không sau một khoảng thời gian (khi quan sát nồng độ khí có giảm thực sự hay không). RL với khái niệm \textit{discounted future reward} xử lý tự nhiên loại tín hiệu trễ này.

  \item \textbf{Cần cân bằng các mục tiêu xung đột}: Phát hiện sớm (sensitivity cao) và tránh cảnh báo sai (specificity cao) là hai mục tiêu thường xung đột. Một hàm reward thiết kế hợp lý cho phép agent học cách cân bằng hai mục tiêu này một cách tự nhiên thay vì phải tinh chỉnh thủ công các ngưỡng.
\end{enumerate}

So với các paradigm khác:
\begin{itemize}
  \item \textbf{Học có giám sát} cần dataset có nhãn cho từng tình huống – điều bất khả thi cho bài toán điều khiển.
  \item \textbf{Học không giám sát} chỉ tìm cấu trúc dữ liệu, không hỗ trợ ra quyết định.
  \item \textbf{Học tăng cường} cho phép học chính sách trực tiếp từ tương tác với bộ mô phỏng, hoàn toàn phù hợp với bài toán đặt ra.
\end{itemize}

\subsection{Các thành phần của một hệ thống RL}

Một hệ thống RL được mô tả bởi các thành phần sau, áp dụng cho bài toán của đề tài:
\begin{itemize}
  \item \textbf{Agent}: Bộ điều khiển khí gas – thực thể ra quyết định trên mỗi mẫu đo lường.
  \item \textbf{Environment}: Bộ mô phỏng vật lý khí gas, có khả năng tiến hóa trạng thái theo các phương trình động lực học.
  \item \textbf{State} $s_t$: Vector mô tả trạng thái hiện tại bao gồm nồng độ khí, nhiệt độ, độ ẩm, và xác suất rủi ro từ LSTM.
  \item \textbf{Action} $a_t$: Một trong bốn hành động rời rạc \{NO\_OP, ALERT, FAN\_ON, CLOSE\_VALVE\}.
  \item \textbf{Reward} $r_t$: Tín hiệu vô hướng được tính từ trạng thái và hành động, thưởng can thiệp hiệu quả và phạt cả bỏ sót lẫn cảnh báo sai.
  \item \textbf{Policy} $\pi(a|s)$: Hàm xác suất ánh xạ từ trạng thái sang hành động, là đối tượng cần học.
  \item \textbf{Value function} $V^\pi(s)$: Ước lượng tổng phần thưởng kỳ vọng từ trạng thái $s$ nếu tuân theo $\pi$.
\end{itemize}

\subsection{Markov Decision Process}

Cơ sở toán học của RL là \textbf{Markov Decision Process} (MDP), được hình thức hóa bởi bộ năm $(\mathcal{S}, \mathcal{A}, \mathcal{P}, \mathcal{R}, \gamma)$:

\begin{itemize}
  \item $\mathcal{S}$: không gian trạng thái (liên tục trong bài toán này).
  \item $\mathcal{A}$: không gian hành động (rời rạc, 4 phần tử).
  \item $\mathcal{P}(s' | s, a)$: phân phối xác suất chuyển trạng thái – mô tả tính ngẫu nhiên nội tại của môi trường.
  \item $\mathcal{R}(s, a)$: hàm phần thưởng.
  \item $\gamma \in [0, 1]$: hệ số chiết khấu, mô tả mức độ ưu tiên phần thưởng tương lai so với hiện tại.
\end{itemize}

Tính chất Markov nghĩa là $\mathcal{P}(s_{t+1} | s_t, a_t, s_{t-1}, a_{t-1}, \ldots) = \mathcal{P}(s_{t+1} | s_t, a_t)$: trạng thái hiện tại đã đủ tổng hợp mọi thông tin cần thiết để dự đoán tương lai. Trong bài toán của đề tài, tính Markov được thỏa mãn xấp xỉ vì xác suất rủi ro $p_{5\min}$ (lấy từ LSTM) đã tổng hợp được lịch sử 60 giây gần nhất vào trạng thái.

Mục tiêu của RL là tìm chính sách tối ưu cực đại hóa giá trị kỳ vọng:
\begin{equation}
  J(\pi) = \mathbb{E}_\pi \left[ \sum_{t=0}^{\infty} \gamma^t r_t \right]
\end{equation}

Hệ số chiết khấu $\gamma$ thường được chọn gần 1 (đề tài sử dụng $\gamma = 0.99$) để khuyến khích agent tính đến cả phần thưởng dài hạn – một yêu cầu quan trọng vì hậu quả của hành động đóng van/bật quạt chỉ thể hiện rõ sau vài giây.

\subsection{Proximal Policy Optimization (PPO)}

Trong họ các thuật toán RL, đề tài lựa chọn \textbf{Proximal Policy Optimization} (PPO) \cite{schulman2017} – một thuật toán policy gradient được nhóm OpenAI đề xuất năm 2017. PPO đã trở thành lựa chọn mặc định trong nhiều ứng dụng RL nhờ sự cân bằng giữa độ ổn định khi huấn luyện và độ phức tạp cài đặt.

Lý do lựa chọn PPO cho bài toán của đề tài:
\begin{itemize}
  \item \textbf{Ổn định khi huấn luyện}: Khác với các thuật toán policy gradient thuần (như REINFORCE) dễ bị cập nhật quá mạnh làm sụp đổ chính sách, PPO sử dụng cơ chế clipping để giới hạn mức thay đổi chính sách trong mỗi bước cập nhật.
  \item \textbf{Hỗ trợ tốt từ thư viện}: Stable-Baselines3 \cite{raffin2021} cung cấp triển khai PPO chất lượng cao, được tối ưu và kiểm thử kỹ lưỡng.
  \item \textbf{Phù hợp với không gian hành động rời rạc}: Bài toán của đề tài có 4 hành động rời rạc, đúng dạng PPO xử lý tốt nhất.
\end{itemize}

\paragraph{Vì sao single-agent PPO chứ không phải multi-agent RL?}
Một câu hỏi thiết kế quan trọng là tại sao chọn kiến trúc single-agent (một policy duy nhất ra quyết định cho toàn bộ thiết bị) thay vì multi-agent RL (MARL) -- vốn là xu hướng nghiên cứu sôi động trong các bài toán IoT phân tán. Đề tài đi theo single-agent với bốn lý do cụ thể:

\begin{enumerate}
  \item \textbf{Bối cảnh triển khai một thiết bị, một van}: Phạm vi của đề tài là phát hiện rò rỉ tại một điểm nguồn (bếp gas hộ gia đình hoặc một khu vực kho nhỏ), với \emph{một} cảm biến và \emph{một} cặp actuator (van + quạt). MARL chỉ có giá trị thực sự khi tồn tại nhiều agent có vùng quyết định giao thoa -- ví dụ nhiều khu vực có thông gió liên thông, hoặc nhiều van phụ thuộc lẫn nhau. Trong bối cảnh hiện tại, không có thứ "phối hợp giữa các agent" để học.

  \item \textbf{Bộ mô phỏng chưa hỗ trợ động lực đa khu vực}: \texttt{World} của simulator (xem §3) là mô hình vật lý đơn-điểm với 4 pha trạng thái, không mô tả các phản ứng vật lý giữa nhiều phòng (luồng khí, thông gió chéo, áp suất). Train MARL trên một môi trường không có tương tác liên-agent sẽ chỉ tái tạo lại single-agent với chi phí mẫu cao hơn nhiều lần. Đây là lý do căn bản (constraint của môi trường), không phải hạn chế thuật toán.

  \item \textbf{Sample efficiency và kích thước policy}: MARL phổ biến (MAPPO, QMIX, MADDPG) thường cần $5\!-\!10$ lần số timestep so với single-agent PPO để hội tụ trên cùng bài toán đơn giản hoá thành single. Với mục tiêu huấn luyện hết trong vài phút CPU, single-agent PPO là lựa chọn hợp lý về mặt tài nguyên cũng như độ phức tạp cài đặt.

  \item \textbf{Non-stationarity và khó debug}: Trong MARL, mỗi agent học trong môi trường mà policy của các agent khác cũng đang thay đổi -- tạo non-stationarity và làm việc gỡ lỗi/đánh giá khó khăn hơn nhiều. Với một báo cáo khoá luận hướng tới \emph{bằng chứng định lượng có thể tái tạo}, single-agent giảm đáng kể số biến đồng thời cần kiểm soát.
\end{enumerate}

Chúng tôi vẫn ghi nhận MARL là hướng mở rộng tự nhiên khi đề tài được scale lên các kịch bản đa-phòng (smart building) hoặc khi muốn phối hợp nhiều van/quạt với mục tiêu toàn cục (xem §5.2 -- Hướng phát triển). Trong khuôn khổ đề tài hiện tại, single-agent PPO là quyết định thiết kế phù hợp với cả \emph{phạm vi bài toán} lẫn \emph{điều kiện tài nguyên} và \emph{tính tái tạo}.

PPO tối ưu hàm mục tiêu clipping:
\begin{equation}
  L^{\text{CLIP}}(\theta) = \mathbb{E}_t \left[ \min\left( r_t(\theta) \hat{A}_t, \; \text{clip}(r_t(\theta), 1-\epsilon, 1+\epsilon) \hat{A}_t \right) \right]
\end{equation}

trong đó:
\begin{itemize}
  \item $r_t(\theta) = \frac{\pi_\theta(a_t | s_t)}{\pi_{\theta_{\text{old}}}(a_t | s_t)}$ là tỷ lệ xác suất giữa chính sách mới và cũ.
  \item $\hat{A}_t$ là ước lượng advantage tại bước $t$.
  \item $\epsilon$ (thường = 0.2) là hệ số clipping giới hạn mức thay đổi.
\end{itemize}

Ý nghĩa trực quan của clipping: khi advantage dương ($\hat{A}_t > 0$, tức hành động tốt hơn kỳ vọng), chúng ta muốn tăng xác suất chọn hành động đó, nhưng không tăng quá $1 + \epsilon$ lần. Ngược lại, khi advantage âm, ta giảm xác suất nhưng không giảm quá $1 - \epsilon$ lần. Cơ chế này ngăn không cho chính sách thay đổi quá nhanh trong một bước cập nhật, đảm bảo tính ổn định.

\subsection{Generalized Advantage Estimation}

Advantage $\hat{A}_t$ trong công thức PPO ước lượng "hành động này tốt hơn trung bình bao nhiêu". PPO thường kết hợp với \textbf{Generalized Advantage Estimation} (GAE) để giảm phương sai trong ước lượng:
\begin{equation}
  \hat{A}_t^{\text{GAE}} = \sum_{l=0}^{\infty} (\gamma \lambda)^l \delta_{t+l}
\end{equation}

trong đó $\delta_t = r_t + \gamma V(s_{t+1}) - V(s_t)$ là Temporal Difference error. Tham số $\lambda \in [0, 1]$ điều chỉnh sự đánh đổi:
\begin{itemize}
  \item $\lambda = 0$: chỉ dùng TD-error một bước, bias cao nhưng variance thấp.
  \item $\lambda = 1$: tương đương Monte Carlo, không bias nhưng variance cao.
  \item Giá trị trung gian (đề tài sử dụng $\lambda = 0.95$) cân bằng tốt giữa bias và variance.
\end{itemize}

\subsection{Thiết kế không gian trạng thái, hành động và phần thưởng}

\textbf{Không gian trạng thái} $\mathcal{S}$ là vector 8 chiều liên tục, mọi thành phần đã được chuẩn hoá về $[0, 1]$:
\begin{itemize}
  \item $gas\_norm = \min(gas / 2000,\,1)$: nồng độ khí hiện tại (ppm).
  \item $temp\_norm = \min(T / 60,\,1)$: nhiệt độ phòng ($^\circ$C).
  \item $hum\_norm = \min(H / 100,\,1)$: độ ẩm tương đối (\%).
  \item $slope\_norm \in [0, 1]$: độ dốc tuyến tính ước lượng từ 30 mẫu cuối, được ánh xạ tuyến tính từ khoảng $[-10, +10]$ ppm/s.
  \item $p_{5\min} \in [0, 1]$: xác suất rủi ro 5 phút từ LSTM forecaster.
  \item $fan\_on \in \{0, 1\}$: trạng thái quạt thông gió.
  \item $valve\_closed \in \{0, 1\}$: trạng thái van gas.
  \item $\Delta t_{\text{act}} = \min(\Delta t / 300,\,1)$: thời gian (giây) kể từ hành động cuối, bão hoà tại 5 phút.
\end{itemize}

Đáng chú ý là việc \textbf{đưa $p_{5\min}$ vào không gian trạng thái} là một quyết định thiết kế quan trọng. Thay vì để RL tự học mọi thứ từ chuỗi đo lường thô (đòi hỏi mô hình lớn và thời gian huấn luyện dài), chúng tôi tận dụng LSTM như một \textit{feature extractor} chuyên biệt cho việc dự báo, và để RL tập trung vào phần ra quyết định. Đây là một dạng \textit{hierarchical learning}: hai mô hình chuyên môn hóa và bổ sung cho nhau. Hai biến trạng thái nội bộ ($fan\_on$, $valve\_closed$) giúp policy biết các can thiệp đã được kích hoạt và tránh lặp lại không cần thiết.

\textbf{Không gian hành động} $\mathcal{A}$ rời rạc với 4 phần tử:
\begin{itemize}
  \item \texttt{0 -- NO\_OP}: không làm gì, chi phí thấp nhất.
  \item \texttt{1 -- ALERT\_USER}: gửi cảnh báo, chi phí trung bình (gây phiền nhiễu nếu sai).
  \item \texttt{2 -- FAN\_ON}: bật quạt, can thiệp vật lý nhẹ giảm khí.
  \item \texttt{3 -- CLOSE\_VALVE}: đóng van, can thiệp mạnh nhất, ngừng nguồn rò rỉ.
\end{itemize}

\textbf{Hàm phần thưởng} được thiết kế theo nguyên tắc \textit{asymmetric cost}: hậu quả của việc bỏ sót rò rỉ (nồng độ vượt ngưỡng nguy hiểm) phải nặng hơn rất nhiều so với cảnh báo sai. Bộ hệ số cụ thể (tham chiếu \texttt{processing/ml/rl/gas\_env.py}) như sau:
\begin{equation}
\begin{aligned}
  r_t = \; & -50\,\mathbf{1}[\text{gas}_t \geq 1000\,\text{ppm}] \\
          & + 0{,}3\,\mathbf{1}[\text{state} \in \{\text{LEAK\_SLOW}, \text{LEAK\_FAST}\} \land \text{gas}_t < 1000] \\
          & - 3\,\mathbf{1}[a_t = 1 \land \text{state} = \text{NORMAL}] \\
          & - 15\,\mathbf{1}[a_t = 2 \land \text{state} = \text{NORMAL}, \text{lần đầu}] \\
          & - 25\,\mathbf{1}[a_t = 3 \land \text{state} = \text{NORMAL}, \text{lần đầu}] \\
          & - 0{,}5\,\mathbf{1}[\text{valve\_closed} \land \text{NORMAL}] - 0{,}1\,\mathbf{1}[\text{fan\_on} \land \text{NORMAL}] \\
          & - 0{,}05 \quad (\text{step cost})
\end{aligned}
\end{equation}

Cấu trúc reward \textit{outcome-based} này phản ánh các quy tắc nghiệp vụ:
\begin{itemize}
  \item \textbf{Phạt rất nặng ($-50$/bước)} khi để khí vượt ngưỡng nguy hiểm: bỏ sót là sai lầm nghiêm trọng nhất, mỗi bước "ngộ độc" đều bị tính.
  \item \textbf{Thưởng theo kết quả ($+0{,}3$/bước)} khi đang rò rỉ mà \emph{giữ được} khí dưới ngưỡng: thưởng cho \emph{kết quả} (kiềm chế thành công) thay vì cho \emph{hành động}. Điều này buộc agent chọn hành động thực sự hiệu quả -- với rò rỉ nhanh (tăng mũ) chỉ có đóng van mới giữ được khí thấp, vì quạt không đủ mạnh.
  \item \textbf{Phạt phân tầng theo chi phí hành động sai trong NORMAL}: $-3$ cho ALERT (chỉ gây phiền), $-15$ cho FAN (tốn điện), $-25$ cho CLOSE\_VALVE (ngắt nguồn cung cấp gas, trải nghiệm xấu).
  \item \textbf{Holding cost} ($-0{,}5$ giữ van, $-0{,}1$ giữ quạt trong NORMAL): tránh việc agent "lười" giữ van đóng vĩnh viễn; kết hợp cơ chế tự khôi phục (mở van/tắt quạt khi phòng an toàn), mỗi sự kiện rò rỉ trở nên độc lập.
  \item \textbf{Step cost $-0{,}05$/bước} để tạo áp lực ra quyết định dứt khoát.
\end{itemize}

Thiết kế này thay thế phần thưởng cố định $+10$ cho mỗi lần bấm FAN/VALVE của phiên bản đầu -- vốn có thể bị \textit{reward hacking} (agent farm điểm bằng cách spam hành động trong lúc rò rỉ mà không thực sự kiềm chế được khí). Việc chuyển sang thưởng theo kết quả là một quyết định thiết kế then chốt, được kiểm chứng định lượng ở Chương~\ref{c:thucnghiemdanhgia}.

\section{Cơ sở dữ liệu thời gian và quan hệ}

\subsection{InfluxDB cho dữ liệu chuỗi thời gian}

InfluxDB là cơ sở dữ liệu chuỗi thời gian mã nguồn mở được tối ưu cho khối lượng ghi cao và truy vấn theo cửa sổ thời gian. Khác với cơ sở dữ liệu quan hệ tổng quát, InfluxDB tổ chức dữ liệu theo mô hình:
\begin{itemize}
  \item \textbf{Measurement}: tương đương "bảng" trong cơ sở dữ liệu quan hệ.
  \item \textbf{Tags}: metadata được đánh chỉ mục, dùng để lọc nhanh (ví dụ \texttt{device\_id}).
  \item \textbf{Fields}: giá trị đo lường thực sự (gas\_ppm, temperature...).
  \item \textbf{Timestamp}: dấu thời gian cho mọi điểm dữ liệu, là khóa chính ngầm định.
\end{itemize}

Việc lựa chọn InfluxDB cho lưu trữ dữ liệu đo lường thay vì PostgreSQL là quyết định có cân nhắc. Một cơ sở dữ liệu quan hệ truyền thống có thể xử lý được nhưng kém hiệu quả khi: (i) ghi với tần suất cao (hàng nghìn điểm/giây), (ii) truy vấn agregation theo cửa sổ thời gian (downsample, last value...), và (iii) lưu trữ dữ liệu lâu dài với retention policy tự động.

Trong hệ thống, measurement \texttt{gas\_readings} được tổ chức:
\begin{itemize}
  \item Tag: \texttt{device\_id}.
  \item Fields: \texttt{gas\_ppm}, \texttt{temperature\_c}, \texttt{humidity\_percent}, \texttt{predicted\_risk\_5min}, \texttt{rl\_action}.
\end{itemize}

\subsection{PostgreSQL cho dữ liệu nghiệp vụ}

Trong khi InfluxDB phù hợp với dữ liệu chuỗi thời gian, các thực thể nghiệp vụ như "sự kiện cảnh báo" và "lịch sử hành động" lại có cấu trúc quan hệ rõ ràng: chúng được truy vấn theo nhiều cách (theo thiết bị, theo mức độ, theo khoảng thời gian, có thể join với metadata...), và đòi hỏi đảm bảo ACID khi cập nhật. PostgreSQL – cơ sở dữ liệu quan hệ mã nguồn mở mạnh mẽ – là lựa chọn tự nhiên.

Trong hệ thống, PostgreSQL lưu trữ hai bảng chính:
\begin{itemize}
  \item \texttt{alert\_events}: lịch sử cảnh báo bao gồm thời điểm, mức độ, nồng độ khí, xác suất rủi ro.
  \item \texttt{action\_events}: nhật ký các hành động RL đã thực thi.
\end{itemize}

Kiến trúc kết hợp hai loại cơ sở dữ liệu (\textit{polyglot persistence}) là một thực hành phổ biến trong các hệ thống dữ liệu hiện đại: chọn đúng công cụ cho đúng loại workload thay vì cố gắng dùng một công cụ duy nhất cho mọi việc.

\section{Telegram Bot API}

Telegram Bot API cho phép phát triển chatbot tương tác với người dùng ngay trong ứng dụng Telegram thông qua HTTP request đến server của Telegram. Việc lựa chọn Telegram làm kênh giao tiếp với người dùng thay vì SMS hoặc ứng dụng di động riêng có ba lý do:

\begin{enumerate}
  \item \textbf{Không yêu cầu phát triển ứng dụng riêng}: Telegram đã có app di động và desktop ổn định, người dùng không cần cài thêm.
  \item \textbf{API hoàn toàn miễn phí}: Khác với SMS gateway thường tính phí trên từng tin nhắn, Telegram cho phép gửi không giới hạn.
  \item \textbf{Hỗ trợ giao diện tương tác phong phú}: Inline keyboards cho phép xây dựng các nút bấm tương tác, ảnh, định dạng văn bản – vượt xa khả năng của SMS.
\end{enumerate}

Trong hệ thống, Telegram bot đóng hai vai trò chính:
\begin{itemize}
  \item \textbf{Kênh cảnh báo đẩy}: Khi bộ điều khiển RL chọn hành động \texttt{ALERT\_USER}, hệ thống tự động gửi message đến người dùng đã đăng ký, kèm thông tin nồng độ khí, xác suất rủi ro và hành động đang thực thi.
  \item \textbf{Giao diện điều khiển từ xa}: Người dùng có thể gửi lệnh đến bot để xem trạng thái hiện tại (\texttt{/status}), lịch sử cảnh báo (\texttt{/history}), hoặc đăng ký/hủy đăng ký nhận thông báo.
\end{itemize}

\section{Docker và Containerization}

Docker là nền tảng containerization cho phép đóng gói ứng dụng cùng toàn bộ phụ thuộc (thư viện, biến môi trường, file cấu hình) vào các \textit{container} nhẹ và cô lập. Khác với máy ảo (Virtual Machine) đóng gói cả hệ điều hành, container chia sẻ kernel với host và chỉ chứa các thành phần đặc thù cho ứng dụng – do đó nhẹ hơn và khởi động nhanh hơn rất nhiều.

Docker Compose là công cụ định nghĩa và quản lý ứng dụng đa container thông qua file YAML khai báo. Đối với đề tài, Docker Compose giải quyết một vấn đề thực tế: hệ thống gồm 10 dịch vụ độc lập (broker, database, processing engine, frontend...) với các phụ thuộc khởi động và cấu hình mạng phức tạp. Việc khởi động thủ công từng dịch vụ một cách đúng thứ tự sẽ rất tốn công và dễ sai sót.

Các dịch vụ trong hệ thống được orchestrate bằng Docker Compose:
\begin{itemize}
  \item \texttt{mosquitto}: MQTT broker.
  \item \texttt{zookeeper} + \texttt{kafka}: Apache Kafka cluster.
  \item \texttt{influxdb}: time-series database.
  \item \texttt{postgresql}: relational database.
  \item \texttt{grafana}: visualization dashboard.
  \item \texttt{mqtt-kafka-bridge}: protocol bridge service.
  \item \texttt{device-simulator}: bộ mô phỏng cảm biến.
  \item \texttt{processing-engine}: Spark + LSTM + RL pipeline.
  \item \texttt{backend}: Node.js API server (kèm Telegram bot).
  \item \texttt{frontend}: Express web dashboard.
\end{itemize}

Đối với một đề tài nghiên cứu, containerization còn mang lại giá trị quan trọng về \textbf{khả năng tái tạo} (reproducibility): toàn bộ môi trường thí nghiệm có thể được khôi phục chính xác trên máy khác chỉ với một lệnh \texttt{docker compose up}. Đây là điều kiện then chốt để kết quả thực nghiệm có thể được kiểm chứng độc lập.
